%% parts/p12/smt-and-exchange-formats.tex --- Part 12.2
\chapter{SMT and Exchange Formats}
\label{sec:p12-smt-and-exchange-formats}

\section{Satisfiability modulo a declared theory}

\begin{hbdefinition}{def:p12-smt-problem}{SMT instance}
Fix a many-sorted first-order signature, a background theory \(T\), and a
declared input fragment.  A \emph{satisfiability modulo theories problem}
asks whether \(T\) together with the existential closure of the input formula
has a model.  The identifier \(\mathsf{SMT}\) is necessary because the
background theory, sorts, and fragment make this different from propositional
SAT; it stands for a parameterized decision problem, not for one logic or one
solver.  Its nonredundancy is witnessed by theory atoms whose consistency is
not determined propositionally.  Relative to the declared signature and
background theory \(T\), a \emph{theory atom} is an atomic object-language
formula over that signature, and a \emph{theory literal} is such an atom or
its negation.  These are interpreted formula components, unlike the fresh
propositional proxy variables used below; the distinction is needed by the
later separation of propositional search from theory checking and by the
later preprocessing that separates nested terms headed by different
component signatures.  A \emph{Boolean
abstraction} is the metalinguistic map that replaces each distinct theory atom
by a propositional proxy, uses the same proxy for repeated occurrences, and
commutes with the Boolean connectives.  It forgets theory interpretation while
preserving propositional structure, so it is neither the original SMT input
nor that later component-signature preprocessing step.  The SMT problem and its Boolean
abstraction are used by the architecture and theory catalog below.
\index{satisfiability modulo theories}\index{SMT}
\index{theory atom}\index{theory literal}\index{Boolean abstraction}
\end{hbdefinition}

\begin{hbdefinition}{def:p12-euf}{EUF}
The quantifier-free theory of \emph{equality with uninterpreted functions}
(\(\mathsf{EUF}\)) has equality and otherwise uninterpreted function symbols;
models may interpret those functions arbitrarily subject to equality
congruence.  On a finite ground input, \emph{congruence closure} is the least
equivalence relation on occurring ground terms that contains the asserted
equalities and is closed under applying the same function symbol to equivalent
argument tuples.  This metalinguistic closure operation is needed to separate
EUF consistency from arithmetic reasoning, is not a synonym for arbitrary
first-order deduction, and is used in Table~\ref{tab:p12-smt-theories}.
\index{equality with uninterpreted functions}\index{congruence closure}
\end{hbdefinition}

\begin{hbmethod}{def:p12-dpllt-architecture}{DPLL(T) architecture}
\index{DPLL(T)}
The name DPLL(\(T\)) is used here only for the following cooperation
architecture, because propositional search and theory consistency have
different responsibilities.  Boolean abstraction replaces theory atoms by
propositional proxies; a SAT engine proposes a Boolean assignment; a
decision procedure for \(T\) checks the selected theory literals; and a theory
conflict returns a \(T\)-valid conflict clause, while theory propagation may
return further entailed literals.  The Boolean engine then learns or branches
until it finds a theory-consistent assignment or refutes the abstraction.
This is a metalinguistic architecture, not a displayed calculus or a
correctness proof; Definition~\ref{def:p12-dpllt-state} and
Theorem~\ref{thm:p12-dpllt-correctness} state the exact transition invariant
and correctness boundary
\cite{NieuwenhuisOliverasTinelli2006}; see also
\cite{BarrettSebastianiSeshiaTinelli2009}.
\end{hbmethod}

\begin{hbdefinition}{def:p12-stably-infinite-convex}{Combination conditions}
A \emph{purification} of a quantifier-free input for disjoint component
signatures repeatedly replaces each maximal subterm headed by a symbol from
the other component signature with a fresh shared variable of the same
declared shared sort and adds the defining equality between that variable and
the replaced subterm to the component whose signature heads that subterm.
Iteration stops when every resulting
literal belongs to one component signature.  The fresh variables make this a
satisfiability-preserving metalinguistic preprocessing step, not Boolean
abstraction: it separates alien term structure rather than replacing whole
theory atoms by propositional proxies.  The Nelson--Oppen statement below uses
the resulting component inputs and their shared variables.

A theory is \emph{stably infinite over its shared sorts} when every satisfiable
quantifier-free input has a model in which each shared sort is infinite.  It is
\emph{convex for shared equalities} when, whenever a conjunction of theory
literals entails a finite disjunction of equalities between shared variables,
it entails one of those equalities.  These properties are needed because
signature disjointness alone does not give the standard Nelson--Oppen
completeness theorem.  For the finite shared variables of each shared sort, an
\emph{equality arrangement} is the complete equality/disequality description
induced by a partition: variables in one block are equated and variables in
different blocks are disequated.  It is compatible with a purified component
input exactly when that input together with the arrangement is satisfiable in
the component theory.  An arrangement is a metalinguistic communication
object, not merely one individually implied equality; that difference is why
nonconvex components require case splitting.  Stable infiniteness and
convexity are semantic properties of component theories, not solver
heuristics: stable infiniteness supports compatible model amalgamation, while
convexity determines whether single equalities suffice or arrangements must
be split.  The theorem and Table~\ref{tab:p12-smt-theories} use all three
distinctions.
\index{purification}\index{stable infiniteness}\index{convex theory}%
\index{equality arrangement}
\end{hbdefinition}

\begin{hbtheorem}{thm:p12-nelson-oppen}{cited=NelsonOppen1979}{Nelson--Oppen combination boundary}
Let quantifier-free inputs be purified into component signatures that are
disjoint apart from equality and explicitly shared sorts.  Suppose each
component theory has a decision procedure for its quantifier-free fragment and
is stably infinite on the shared sorts.  Then satisfiability in their union is
decidable by exchanging compatible equality arrangements on the shared
variables.  For convex components, communication of individually implied
equalities suffices; nonconvex components require disjunctive case splitting.
Disjointness supplies sound separation, but without the stated
stable-infiniteness hypothesis this completeness conclusion does not follow.
\end{hbtheorem}

\begin{landscape}
\thispagestyle{plain}
\begingroup\scriptsize
\setlength{\tabcolsep}{3pt}
\begin{longtable}{@{}p{.95in}p{1.90in}p{2.15in}p{2.15in}@{}}
\caption{Named SMT theories and combination boundaries.}
\label{tab:p12-smt-theories}\\
\toprule
Theory/fragment & Exact source-defined problem & Decidability and complexity boundary
& Nelson--Oppen field \\ \midrule
\endfirsthead
\caption[]{Named SMT theories and combination boundaries (continued)}\\
\toprule
Theory/fragment & Exact source-defined problem & Decidability and complexity boundary
& Nelson--Oppen field \\ \midrule
\endhead
\midrule\multicolumn{4}{r}{\emph{continued on next page}}\\
\endfoot
\endlastfoot
\(\mathsf{EUF}\) & Quantifier-free equality with uninterpreted functions &
Decidable by congruence closure for conjunctions; arbitrary Boolean structure
contains propositional SAT and is NP-complete under the usual finite-term
encoding & Stably infinite and convex; a standard component
\cite{BarrettSebastianiSeshiaTinelli2009} \\ \addlinespace
\texttt{ArraysEx} / \texttt{QF\_AX} & Closed quantifier-free formulas over
an arbitrary expansion of the SMT-LIB 2.7 \texttt{ArraysEx} theory with free
sort and constant symbols; \texttt{ArraysEx} supplies \texttt{Array X Y},
\texttt{select}, \texttt{store}, both read-over-write laws, and extensionality &
Decidable; arbitrary Boolean
structure is NP-complete in the standard term encoding.  Adding further array
operations changes the claim & The displayed extensional functional-array
theory, without a constant-array extension, is stably infinite and nonconvex.
Under the theorem's disjoint-signature/shared-sort hypotheses and
stable-infiniteness of the other component, complete cooperation therefore
uses compatible equality arrangements and case splits rather than only
individually implied equalities.  No blanket combination claim is made for
further expansions
\cite{SMTLIB2025,BarrettSebastianiSeshiaTinelli2009} \\
\addlinespace
\texttt{FixedSizeBitVectors} / \texttt{QF\_BV} & Closed quantifier-free
formulas over arbitrary expansions of the SMT-LIB 2.7
\texttt{FixedSizeBitVectors} signature, with the additional operators and
abbreviations declared by \texttt{QF\_BV} and free constants of sorts
\texttt{(\_ BitVec m)} for numerals \(m>0\) &
Finite and decidable.  Replacing each bit-vector operation by its explicit
propositional circuit reduces declared fixed widths to SAT;
compact binary-coded widths can cause exponential expansion, so a complexity
class must state its width encoding & Not stably infinite at fixed width;
ordinary Nelson--Oppen is not directly applicable \cite{SMTLIB2025} \\
\addlinespace
Linear rational/real arithmetic & Quantifier-free linear constraints over
\(\mathbb Q\) or \(\mathbb R\) & Conjunctive feasibility is polynomial-time;
arbitrary Boolean combinations are NP-complete because SAT embeds &
Stably infinite and convex in the standard signatures
\cite{DutertreDeMoura2006} \\ \addlinespace
Linear integer arithmetic & Quantifier-free Presburger constraints with
binary-coded coefficients & Decidable and NP-complete for the displayed
quantifier-free input form & Stably infinite but nonconvex; complete
combination requires equality arrangements or case splitting \\ \bottomrule
\end{longtable}
\endgroup
\end{landscape}

\section{SMT-LIB and TPTP are interfaces}

\begin{hbdefinition}{def:p12-smtlib-interface}{SMT-LIB interface boundary}
SMT-LIB 2.7 is a versioned command language and family of theory/logic
declarations for communicating SMT problems and responses.  It is not itself
a calculus or one semantic theory.  A reproducible SMT-LIB record therefore
fixes the logic declaration, sort and symbol declarations, assertion stack,
options, resource limits, solver/version, and the exact response token
\texttt{sat}, \texttt{unsat}, or \texttt{unknown}; a model or proof artifact
must be requested and checked separately.  This interface name is necessary
to distinguish transport syntax from Definition~\ref{def:p12-smt-problem},
and that distinction is used by the automated-result audit
\cite{SMTLIB2025}.
\index{SMT-LIB}
\end{hbdefinition}

\begin{hbdefinition}{def:p12-tptp-interface}{TPTP interface boundary}
TPTP's BNF revision 9.3.0.3 supplies concrete languages and annotated-formula
roles for exchanging theorem-proving problems.  It is not a single logic or a
proof of any submitted conjecture.  A reproducible record fixes the grammar
revision, declared TPTP language, include files, formula roles, conjecture polarity,
tool/version, resource limits, returned status, and any independently checked
proof or model.  The name is needed because roles such as axiom and conjecture
affect problem meaning rather than merely typography, and the distinction is
used in the release audit \cite{TPTP}.
\index{TPTP}
\end{hbdefinition}

\begin{longtable}{@{}p{.18\textwidth}p{.34\textwidth}p{.38\textwidth}@{}}
\caption{Exchange-format audit.}
\label{tab:p12-exchange-format-audit}\\
\toprule
Layer & Record & Failure prevented \\ \midrule
\endfirsthead
\caption[]{Exchange-format audit (continued)}\\
\toprule
Layer & Record & Failure prevented \\ \midrule
\endhead
\midrule\multicolumn{3}{r}{\emph{continued on next page}}\\
\endfoot
\endlastfoot
Language & SMT-LIB logic/theory declarations or TPTP language and roles &
Treating a parser choice as a semantic theorem \\ \addlinespace
Environment & Sorts, symbols, imported assertions, include files, and scopes &
Running a stronger or different problem than the reported one \\ \addlinespace
Execution & Tool build, options, limits, random seed where relevant, and exact
response & Reading timeout or \texttt{unknown} as a negative answer \\ \addlinespace
Evidence & Model, proof log, checker/version, and checker result & Treating an
unverified solver response as an independently checked certificate \\
\bottomrule
\end{longtable}

\section{DPLL(\(T\)) transitions and invariants}

\begin{hbdefinition}{def:p12-dpllt-state}{DPLL(\(T\)) state}
For a fixed Boolean abstraction \(B(F)\), a state is
\((D,\tau,\mathcal A)\): \(D\) is the propositional clause database, initially
the clausified abstraction; \(\tau\) is a decision-level trail of signed proxy
literals with Boolean antecedents; and \(\mathcal A\) is polarity respecting:
if \(p\) represents theory atom \(L\), then
\(\mathcal A(p)=L\) and \(\mathcal A(\neg p)=\neg L\). The theory checker
receives \(\mathcal A[\tau]\). The transitions are Boolean decide/propagate,
theory check, theory propagate, conflict-clause learning, Boolean backjump,
and success when the total trail is propositionally and \(T\)-consistent.
The state name is needed because Boolean implication and theory implication
have distinct evidence. It stands as an operational state, differs from plain
CDCL by \(\mathcal A\) and the theory transitions, and is used in the
correctness invariant \cite{NieuwenhuisOliverasTinelli2006}.
\index{DPLL(T)!transition state}
\end{hbdefinition}

A theory propagation \(\mathcal A[\tau]\models_T L\) returns the valid clause
\(\neg\bigwedge\mathcal A[\tau]\lor L\); a theory conflict
\(\mathcal A[\tau]\models_T\bot\) returns a clause consisting of the
complements of a finite inconsistent core. After Boolean abstraction these are
ordinary learned clauses. Their invariant is not that the SAT engine derived
them propositionally from the old database, but that their concretizations are
valid in every \(T\)-model.

\begin{hbtheorem}{thm:p12-dpllt-correctness}{cited=NieuwenhuisOliverasTinelli2006}{DPLL(\(T\)) correctness boundary}
Assume clausification is a definitional extension: every original model
extends to a model of the clauses and every clause model projects, by
forgetting fresh symbols, to a model of the original formula. Assume every theory propagation and
conflict clause is \(T\)-valid, and the theory checker is sound and complete for
the declared finite literal fragment. Then an accepted total state yields a
model of the original SMT instance, and a derivation of the empty clause proves
it \(T\)-unsatisfiable. If the Boolean search is fair over finitely many proxies
and the theory procedure terminates on every finite trail, the transition
system decides the fragment.
\end{hbtheorem}

\noindent\emph{Proof.}
Inductively, every clause in \(D\) is valid under \(T\) whenever the original
input is: initial clauses come from the stated definitional clausification,
Boolean learned clauses are resolution consequences, and theory clauses are
valid by hypothesis. Thus the empty clause is a sound refutation. At success,
the Boolean assignment satisfies the abstraction and \(\mathcal A[\tau]\) is
jointly \(T\)-satisfiable; a witnessing extended \(T\)-model projects by the
stated invariant to a model of the original Boolean combination. Finiteness and fairness exhaust the Boolean
assignments unless one succeeds, while terminating complete theory checks
classify each selected conjunction. Dropping any hypothesis invalidates the
corresponding conclusion.

For ground EUF, congruence closure maintains equivalence classes with a
signature table keyed by function symbol and argument classes; a disequality
whose endpoints merge is a conflict. Explanation edges recover a finite subset
of asserted equalities justifying each merge. Linear rational arithmetic may
use exact simplex pivots and return a Farkas-style infeasibility combination;
linear integer arithmetic additionally needs integrality reasoning such as
branching or cutting planes. Floating-point arithmetic, nonlinear arithmetic,
arrays, bit-vectors, and strings require their own declared procedures and
proof objects; none inherits the rational-linear proof automatically
\cite{DutertreDeMoura2006,BarrettSebastianiSeshiaTinelli2009}.

\noindent\emph{Nelson--Oppen proof boundary.}
Soundness follows because purification is definitional and each component sees
only equalities shared by both. For completeness, choose a compatible complete
arrangement of the finitely many shared variables. Each component has a model
of its purified literals plus the arrangement; stable infiniteness permits the
shared equivalence classes to be embedded into equally large shared-sort
domains, after which the disjoint nonshared signatures are amalgamated. Convex
theories can discover the arrangement by communicating single implied
equalities; nonconvex theories must split arrangements. This establishes the
theorem under its exact disjointness and stable-infiniteness hypotheses
\cite{NelsonOppen1979}.
